\chapter{ทัศนศาสตร์ กัมมันตภาพรังสี และการประยุกต์}
\label{ch:ch13_optics_radioactivity}
\index{ทัศนศาสตร์และกัมมันตรังสี}

\begin{conceptbox}[title=วัตถุประสงค์การเรียนรู้ประจำบท]
เมื่อสิ้นสุดการศึกษาบทเรียนนี้ นักศึกษาสามารถ
\begin{enumerate}
    \item เมื่อศึกษาบทนี้จบแล้ว ผู้เรียนสามารถ
    \item อธิบายกฎการสะท้อนของแสงและลักษณะการเกิดภาพจากกระจกเงาราบและกระจกเงาโค้ง (เว้าและนูน)
    \item อธิบายปรากฏการณ์การหักเหของแสง และประยุกต์ใช้กฎของสเนลล์ในการคำนวณมุมหักเห
    \item อธิบายเงื่อนไขการเกิดการสะท้อนกลับหมด (total internal reflection) และคำนวณมุมวิกฤตได้
    \item ใช้สมการเลนส์บางและสูตรกำลังขยายในการหาตำแหน่ง ชนิด และขนาดของภาพที่เกิดจากเลนส์
\end{enumerate}
\end{conceptbox}

\section{ทฤษฎีและหลักการพื้นฐาน}

\begin{bioappbox}
\textbf{ทัศนศาสตร์ของกล้องจุลทรรศน์และการใช้รังสีในงานชีววิทยา}

กล้องจุลทรรศน์แบบใช้แสง (Compound Light Microscope) อาศัยการหักเหของแสงผ่านเลนส์ใกล้วัตถุและเลนส์ใกล้ตาเพื่อให้ได้กำลังขยายรวม $M = M_o \times M_e$ ขีดจำกัดในการแยกชัดของกล้อง (Limit of Resolution) ถูกกำหนดโดยเกณฑ์ของเรย์ลีห์ $d = \frac{0.61\lambda}{NA}$ ในขณะที่ปรากฏการณ์กัมมันตภาพรังสี เช่น รังสีแกมมา ถูกนำมาใช้ในการกำจัดเชื้อจุลินทรีย์ในเวชภัณฑ์และผลผลิตทางการเกษตร รวมถึงการใช้ไอโซโทปรังสีคาร์บอน-14 ในการหาอายุซากดึกดำบรรพ์และติดตามเส้นทางเมแทบอลิซึม
\end{bioappbox}

\begin{labbox}
1. ตรวจสอบตำแหน่งศูนย์ (Zero Error) ของเครื่องมือวัดทุกชนิดก่อนเริ่มบันทึกข้อมูล
2. ทำการวัดซ้ำอย่างน้อย 3 ถึง 5 ครั้งในแต่ละจุดทดลองเพื่อคำนวณหาค่าเฉลี่ยและค่าเบี่ยงเบนมาตรฐาน
3. บันทึกผลการวัดพร้อมระบุหน่วยเอสไอ (SI Units) และค่าความคลาดเคลื่อนของการวัดเสมอ
\end{labbox}

\section{ตัวอย่างโจทย์และการวิเคราะห์เชิงฟิสิกส์}

\begin{examplebox}
บทนี้เป็นบทสุดท้ายของเอกสารประกอบการสอนวิชาฟิสิกส์พื้นฐาน โดยนำเสนอเนื้อหาสองส่วนที่แม้จะดูแตกต่างกันในระดับปรากฏการณ์ คือ ทัศนศาสตร์ซึ่งว่าด้วยพฤติกรรมของแสงในระดับมหภาค และฟิสิกส์นิวเคลียร์ซึ่งว่าด้วยปรากฏการณ์ในระดับจุลภาค แต่ทั้งสองส่วนต่างก็เป็นตัวอย่างที่แสดงให้เห็นว่าหลักการฟิสิกส์พื้นฐานที่ได้ศึกษามาตลอดทั้งเล่ม ตั้งแต่กลศาสตร์ คลื่น ไฟฟ้า และแม่เหล็ก สามารถนำมาประยุกต์อธิบายปรากฏการณ์ธรรมชาติที่หลากหลายและเชื่อมโยงเข้ากับเทคโนโลยีสมัยใหม่ได้อย่างเป็นระบบ

เนื้อหาหลัก

13.1 การสะท้อนของแสง

การสะท้อนของแสง (reflection) คือปรากฏการณ์ที่แสงเปลี่ยนทิศทางการเดินทางเมื่อกระทบผิวของตัวกลางแล้วสะท้อนกลับเข้าสู่ตัวกลางเดิม กฎการสะท้อนกล่าวว่า มุมตกกระทบ (angle of incidence) ซึ่งวัดจากเส้นแนวฉาก (normal) ที่ตั้งฉากกับผิวสะท้อน ณ จุดตกกระทบ มีค่าเท่ากับมุมสะท้อน (angle of reflection) เสมอ และรังสีตกกระทบ เส้นแนวฉาก และรังสีสะท้อน ทั้งสามอยู่ในระนาบเดียวกัน กล่าวโดยสัญลักษณ์คือ $\theta$ตกกระทบ = $\theta$สะท้อน
\end{examplebox}

\section{แบบฝึกหัดทบทวนและคำถามท้ายบท}

\begin{enumerate}
    \item อธิบายกฎการสะท้อนของแสง และลักษณะภาพที่เกิดจากกระจกเงาราบ
    \item แสงเดินทางจากอากาศ (n = 1.00) เข้าสู่แก้ว (n = 1.50) ทำมุมตกกระทบ 40$^\circ$ จงหามุมหักเห
    \item อธิบายความแตกต่างของลักษณะภาพที่เกิดจากกระจกเงาเว้าและกระจกเงานูน พร้อมยกตัวอย่างการใช้งาน
    \item อธิบายเงื่อนไขการเกิดการสะท้อนกลับหมด (TIR) และประโยชน์ของปรากฏการณ์นี้ในเส้นใยนำแสง
    \item น้ำมีดัชนีหักเหเท่ากับ 1.33 จงหามุมวิกฤตของแสงที่เดินทางจากน้ำสู่อากาศ
    \item วัตถุวางห่างจากเลนส์นูนที่มีความยาวโฟกัส 15 cm เป็นระยะ 45 cm จงหาตำแหน่งภาพและกำลังขยาย
\end{enumerate}

% สิ้นสุดบทเรียน
